\section*{الفصل \ref{ch:g5:people-and-nature} --- الناس والطبيعة}

\begin{solution}{exo:g5:people-and-nature:1}
الحقول، وتُدار لمحصول واحد مختار؛ \omterm{prop:g5:people-and-nature:managed}{والغابات}، وتُدار للأخشاب؛
والمروج، وتُبقى مفتوحة للدريس أو للرعي؛ والمدن والطرق والسدود،
وتُبنى للسكن والتنقل. (وأي أربعة منها مع أغراضها.)
\end{solution}

\begin{solution}{exo:g5:people-and-nature:2}
إن تُرك بلا جزّ ولا رعي عاد أدغالًا في سنوات قليلة، ثم حرشًا ---
فيختفي \omterm{def:g2:where-animals-live:habitat}{موئل} المرج المفتوح المزهر.
\end{solution}

\begin{solution}{exo:g5:people-and-nature:3}
لا تأخذ أسرع مما يتجدد؛ وأبقِ العمّال (\omterm{def:g5:biodiversity:species}{الأنواع} التي يعتمد
الاستعمال على خدماتها المجانية)؛ وردّ ما تستعير (أغلق الحلقات، ولا
تترك فضلات تتكوّم).
\end{solution}

\begin{solution}{exo:g5:people-and-nature:4}
الأسيجة مقتلعة في مقابل محفوظة: فالسفح الأول يفقد ماسكات التربة
وبيوت آكلي الآفات. والتربة العارية شتاءً في مقابل المغطاة: فتربة
الأول تُغسل إلى أسفل في الأمطار. والمحصول الأوحد العملاق في مقابل
التناوب: فالأول يدعو أي آفة لذلك المحصول إلى اكتساح السفح كله، ولا
يريح الأرض أبدًا.
\end{solution}

\begin{solution}{exo:g5:people-and-nature:5}
أحصِ ما تنميه الغابة في السنة؛ ولا تقطع أكثر منه. وهي القاعدة
الأولى --- لا تأخذ أسرع مما يتجدد --- وتغلّ الغابة خشبًا كل سنة بلا
نهاية، وتبقى غابة حية.
\end{solution}

\begin{solution}{exo:g5:people-and-nature:6}
مثلًا: السمامة تعشّش على ``جروف'' الخرسانة؛ والثعالب تجوب
الحاويات والحدائق؛ \omterm{def:g1:plants-around-us:plant}{والنباتات} البرية تتجذّر في شقوق الإسفلت
\omterm{def:g1:my-body:joint}{ومفاصل} الجدران (وكذلك الحمام على الأفاريز والصقور على الأبراج).
\end{solution}

\begin{solution}{exo:g5:people-and-nature:7}
للناس: ماء أبطأ وتخزين للفيضان في المروج الرطبة المستعادة ---
وفيضان أقل في الأسفل. وللشبكة: عادت \omterm{def:g2:where-animals-live:habitat}{موائل} الانحناءات، وعادت معها
\omterm{prop:g3:sorting-living-things:groups}{الأسماك} والرعّاشات ومروج صيد اللقالق.
\end{solution}

\begin{solution}{exo:g5:people-and-nature:8}
الجواب مفتوح. ومما يُوجد عادة: أكتاف طرق مجزوزة، وأشجار مقلَّمة،
وأحواض مغروسة، وأرصفة منزوعة الأعشاب --- وربما قمة جدار أو خندق
واحد غير مُدار، وهو في الغالب أبرّ بقعة في المنظر.
\end{solution}

\begin{solution}{exo:g5:people-and-nature:9}
ألّا نمسّها البتة سيدمر بعض \omterm{def:g2:where-animals-live:habitat}{الموائل} فعلًا: فالمرج المزهر موجود
\emph{بسبب} الجزّ، وينغلق من غيره. فحماية الطبيعة ليست انعدام
المسّ بل حسن الإدارة --- والسؤال الصادق هو كيف نمسّ: أي القواعد
يحترمها الاستعمال، لا أوُضعت الأيدي أصلًا أم لا.
\end{solution}

\begin{solution}{exo:g5:people-and-nature:10}
القاعدة الثانية مخروقة صراحةً (فقد طُرد العمّال مع الأسيجة)،
والثالثة جزئيًّا (فالتربة المغسولة لا تُردّ أبدًا)، والأولى في خطر
(فالتربة تُفقد أسرع مما تتكوّن). وأكثر تغيير واحد إصلاحًا: إعادة
غرس الأسيجة --- فهي تمسك التربة وتعيد إسكان آكلي الآفات في ضربة
واحدة.
\end{solution}

\begin{solution}{exo:g5:people-and-nature:11}
الصيد المضاعف أخذ السمك أسرع بكثير مما تتكاثر به الأسراب ---
فأنفق مخزون التكاثر نفسه. وإنفاق رأس المال يبدو باهرًا دائمًا ما
دام باقيًا: فقد امتلأت العنابر لأن سمك السنة القادمة كان يُنزَل
هذه السنة بالضبط. أما الحرّاج فيحصد نمو السنة وحده ويُبقي المخزون
النامي كاملًا --- الحساب نفسه بإشارة مقابلة، ولهذا تغلّ الغابة إلى
الأبد ووقف الأسطول في السنة السادسة.
\end{solution}
